\section{与Sobolev空间的联系}
\label{sec:besov-sobolev}

Besov空间是为补充Sobolev条件而引入的，自然要问二者在重合处是否一致。本节只对照第~\ref{chap:sobolev-spaces}章已经定义的一阶平方可积弱导数：在单位周期设定下，$B^1_{2,2}$恰好由端点值相同的$H^1(0,1)$函数组成，范数等价。随后说明端点条件与二阶差分各自为何必要，再用其他参数说明Besov空间在Sobolev空间之外增加了什么。

\paragraph{定理与证明}
记$H^1_{\mathrm{per}}=\{f\in H^1(0,1):\tilde f(0)=\tilde f(1)\}$，其中$\tilde f$是$f$的绝对连续代表，这类函数按周期$1$延拓后在接缝处连续。

\begin{theorem}[$B^1_{2,2}$与周期$H^1$]
\label{thm:besov-b122-h1}
在单位周期上，$B^1_{2,2}(\mathbb T)=H^1_{\mathrm{per}}$，并且存在常数$C$，使$C^{-1}\lVert f\rVert_{H^1(0,1)}\leq\lVert f\rVert_{B^1_{2,2}}\leq C\lVert f\rVert_{H^1(0,1)}$。
\end{theorem}

证明借助周期Fourier系数$\widehat f(k)=\int_0^1f(x)e^{-2\pi\mathrm ikx}\,\mathrm dx$，分两步：先说明$H^1_{\mathrm{per}}$恰好由$\sum_kk^2|\widehat f(k)|^2<\infty$的函数组成，再说明$B^1_{2,2}$的半范数平方与这个和等价。

\emph{第一步：Fourier系数刻画$H^1_{\mathrm{per}}$。}设$f\in H^1_{\mathrm{per}}$。对绝对连续代表分部积分，
\begin{equation*}
\begin{aligned}
\widehat{f'}(k)&=\Bigl[\tilde f(x)e^{-2\pi\mathrm ikx}\Bigr]_0^1+2\pi\mathrm ik\int_0^1\tilde f(x)e^{-2\pi\mathrm ikx}\,\mathrm dx\\
&=\bigl(\tilde f(1)-\tilde f(0)\bigr)+2\pi\mathrm ik\,\widehat f(k)=2\pi\mathrm ik\,\widehat f(k),
\end{aligned}
\end{equation*}
端点条件正好消去边界项。由Parseval等式，$\lVert f'\rVert_2^2=4\pi^2\sum_kk^2|\widehat f(k)|^2<\infty$。反过来，设$\sum_kk^2|\widehat f(k)|^2<\infty$，则$g=\sum_k2\pi\mathrm ik\widehat f(k)e^{2\pi\mathrm ikx}\in L^2$，且$\widehat g(0)=0$。令$F(x)=c+\int_0^xg$，常数$c$使$\int_0^1F=\widehat f(0)$。$F$绝对连续，$F(1)-F(0)=\widehat g(0)=0$；对$k\neq0$分部积分得$\widehat F(k)=\widehat g(k)/(2\pi\mathrm ik)=\widehat f(k)$。Fourier系数全部相同，由完备性$F=f$几乎处处，所以$f\in H^1_{\mathrm{per}}$。

\emph{第二步：半范数与系数和等价。}$s=1$时须取$r\geq2$，这里取$r=2$。由式\eqref{eq:besov-difference-fourier}，$\lVert\Delta_h^2f\rVert_2^2=\sum_k(2|\sin(\pi kh)|)^4|\widehat f(k)|^2$。下界：$\omega_2(f,t)_2\geq\lVert\Delta_t^2f\rVert_2$，交换求和与积分，
\[
|f|_{B^1_{2,2}}^2\geq\sum_k|\widehat f(k)|^2\int_0^1t^{-3}\,16\sin^4(\pi kt)\,\mathrm dt .
\]
对$k\neq0$换元$u=|k|t$，内层积分等于$k^2\int_0^{|k|}16u^{-3}\sin^4(\pi u)\,\mathrm du\geq c_1k^2$，其中$c_1=16\int_0^1u^{-3}\sin^4(\pi u)\,\mathrm du$为正的有限数（被积函数在$u\to0$时约为$\pi^4u$）。上界：$(2|\sin(\pi kh)|)^4\leq\min\bigl(16,(2\pi|k||h|)^4\bigr)$，右端随$|h|$不减，所以$\omega_2(f,t)_2^2\leq\sum_k\min(16,(2\pi|k|t)^4)|\widehat f(k)|^2$；同样换元，$\int_0^1t^{-3}\min(16,(2\pi|k|t)^4)\,\mathrm dt\leq C_1k^2$，$C_1=\int_0^\infty u^{-3}\min(16,(2\pi u)^4)\,\mathrm du<\infty$。于是
\[
c_1\sum_kk^2|\widehat f(k)|^2\leq|f|_{B^1_{2,2}}^2\leq C_1\sum_kk^2|\widehat f(k)|^2 .
\]
两步合并，$f\in B^1_{2,2}$当且仅当$f\in H^1_{\mathrm{per}}$，且$|f|_{B^1_{2,2}}$与$\lVert f'\rVert_2$等价；再加上$\lVert f\rVert_2$，两个范数等价。

\paragraph{端点条件的作用}
不能把全部$H^1(0,1)$当作周期空间。$f(x)=x$属于$H^1(0,1)$，但$\tilde f(0)=0\neq1=\tilde f(1)$，按周期延拓后在每个整数点跳跃。直接计算，$k\neq0$时$\widehat f(k)=-1/(2\pi\mathrm ik)$，于是$k^2|\widehat f(k)|^2=1/(4\pi^2)$不可和，$f\notin B^1_{2,2}$。它的系数与周期阶跃同样按$1/|k|$衰减，周期延拓后的光滑阶数只有$1/2$。第一步的分部积分清楚地显示了原因：边界项$\tilde f(1)-\tilde f(0)$不为零时，$\widehat{f'}(k)$与$2\pi\mathrm ik\widehat f(k)$相差一个常数，这个常数正是接缝处的跳跃。

\paragraph{二阶差分的必要性}
$s=1$时差分阶数必须取$r\geq2$。若用一阶差分，由第~\ref{sec:besov-modulus}小节的饱和现象，任何非常数的$f$都有某个$\widehat f(k)\neq0$，$k\neq0$；当$|k|t\leq1/2$时，$\omega_1(f,t)_2\geq2|\sin(\pi kt)||\widehat f(k)|\geq4|k|t|\widehat f(k)|$，所以$t^{-1}\omega_1(f,t)_2$在小$t$处有正的下界，$\int_0(t^{-1}\omega_1)^2\,\mathrm dt/t$发散。用一阶差分写出的“$B^1_{2,2}$”只含常数。小波刻画也有同样的要求：所用小波须有超过一阶的正则性与消失矩，Haar小波不能用于$s=1$。

\paragraph{其他参数}
同样的Fourier计算适用于其他光滑阶数。$0<s<1$时取$r=1$，$\int_0^1t^{-1-2s}4\sin^2(\pi kt)\,\mathrm dt$与$|k|^{2s}$同阶，所以$B^s_{2,2}$由$\sum_k(1+k^2)^s|\widehat f(k)|^2<\infty$刻画，它是分数阶的平方可积光滑性，整数阶的$H^1_{\mathrm{per}}$是其中$s=1$的情形。可以证明，对正整数$s$，$B^s_{2,2}$与$s$阶的周期$H^s$一致，这里不给出证明。

$p\neq2$时情况不同。对$1<p<\infty$，可以证明周期的$W^{1,p}$介于$B^1_{p,\min(p,2)}$与$B^1_{p,\max(p,2)}$之间，一般不等于任何一个$B^1_{p,q}$；整数阶Sobolev空间与Besov空间只在$p=2$时简单重合。这一结论这里不给出证明。

Besov空间在Sobolev空间之外增加的，正是本章开头要求的两项描述。光滑阶数$s$可以取分数，于是没有弱导数的周期阶跃也有确切位置：它属于$B^{1/2}_{2,\infty}$，并属于每个$s<1/2$的$B^s_{2,q}$。汇总指数$q$可以选择，于是只在统一速率上相同、跨尺度分布不同的函数，例如第~\ref{sec:besov-distribution-aggregation}小节的两种权重，或例~\ref{ex:besov-one-coefficient-per-scale}，可以被区分开。积分指数$p$则与$s$配合，决定像$B^s_{\tau,\tau}$那样允许稀疏大系数的空间，这对应于最佳$N$项逼近。
